\lysh{38810}{4}

\begin{alist}
\item Έστω $x$ το κεφάλαιο που θα επενδύσει ο Ιάσονας και $y = x + 4000$ της Νεφέλης. Η συνολική επένδυση των δύο φίλων θα είναι το άθροισμα των δύο κεφαλαίων. Επομένως $f(x) = x + (x + 4000) = 2x + 4000$.
Άρα η συνάρτηση που εκφράζει το συνολικό ποσό που θα επενδύσουν η Νεφέλη και ο Ιάσονας ως συνάρτηση του $x$, δηλαδή του κεφαλαίου του Ιάσωνα είναι η:
$f(x) = 2x + 4000$.

\item
\begin{rlist}
\item $16000 \le f(x) \le 30000 \Leftrightarrow 16000 \le 2x + 4000 \le 30000 \Leftrightarrow$
\[16000 - 4000 \le 2x \le 30000 - 4000 \Leftrightarrow 12000 \le 2x \le 26000 \Leftrightarrow\]
$6000 \le x \le 13000$ (1)
Άρα ο Ιάσονας μπορεί να επενδύσει από $6000$\euro\ έως και $13000$\euro.

\item Αν η Νεφέλη επενδύσει $20000$\euro, ο Ιάσονας θα πρέπει να επενδύσει:
$20000 - 4000 = 16000$\euro.
Όμως το $16000$ είναι μεγαλύτερο από το $13000$, που είναι το μέγιστο ποσό που μπορεί να επενδύσει ο Ιάσονας, για να συνεχίσει να είναι ασφαλής η επένδυση. Άρα η Νεφέλη δεν μπορεί να επενδύσει όλο το ποσό των $20000$\euro.

\item Αφού η Νεφέλη επένδυσε $4000$\euro\ περισσότερα από τον Ιάσονα, για να βρούμε το διάστημα των δυνατών ποσών που μπορεί να επενδύσει, θα πάρουμε την ανίσωση (1) του $\beta)$ ερωτήματος και θα προσθέσουμε $4000$ σε όλα τα μέλη:
$6000 + 4000 \le x + 4000 \le 13000 + 4000 \Leftrightarrow 10000 \le y \le 17000$ (2).
Άρα η Νεφέλη μπορεί να επενδύσει από $10000$\euro\ έως και $17000$\euro.
\end{rlist}

\item Έστω ότι $\alpha$ είναι το ποσό που έχει αποταμιεύσει ο Ιάσονας. Η Νεφέλη επενδύει $17000$\euro\ και δανείζει $3000$\euro\ στον Ιάσονα. Άρα το συνολικό ποσό που επενδύουν μαζί είναι:
\[K = 17000 + (\alpha + 3000) = 20000 + \alpha.\]
Το ποσό αυτό όμως δεν πρέπει να υπερβαίνει τις $30000$\euro, δηλαδή:
\[20000 + \alpha \le 30000 \Leftrightarrow \alpha \le 30000 - 20000 \Leftrightarrow \alpha \le 10000.\]
Άρα ο Ιάσονας πρέπει να έχει αποταμιεύσει το πολύ $10000$\euro, ώστε το συνολικό ποσό επένδυσης της Νεφέλης και του Ιάσονα μαζί να μην υπερβαίνει τις $30000$\euro.
\end{alist}
